\section{Small loci, selected curves, and the actual sharp ideal}
\label{sec:entry-nice-sharp}
We continue with the single entry algebra and its
completed charts in Section~\ref{sec:entry-space-flat}.
All dimensions and cohomology constants below
belong to the arithmetic frame already constructed.

\subsection{The two small loci}
The lower-triangular characteristic fibre has
dimension at most \(2C_0\).
Here is a generator proof, valid over rings.
In rank three its matrices are determined by
the lower functions \((y,a)\) used in the
explicit interpolation.
The linear space of possible \(y\)'s injects
into the strict \(H^1\) for \(\Psi^{-1}\)
over the CM subgroup.
The marked equality at \(g_0,g_0^*\) kills a
possible coboundary: it gives
\((\Psi(g_0)-\Psi(g_0)^{-1})h=0\).
Thus its dimension is at most \(C_0\).
The space of strict linear differences of
the \(a\)'s has dimension at most \(C_0\),
by normalization at \(c\) and the strict
\(H^1\) for \(\Psi^{-2}\delta\).

Choose a basis of the \(y\)-space.
Finitely many evaluations recover its
coefficients, so all values of \(y\) belong
to an algebra generated by at most \(C_0\)
elements.
Modulo this algebra, the quadratic term
\(y(g^*)y(h)\) in the equation for \(a\)
vanishes.
Its local primitive products vanish there
too: separated local characters recover
each primitive from one \(y\)-value.
The resulting additive quotient function
is in the strict linear \(a\)-space.
Choose representatives for its at most
\(C_0\) coefficients.
Together these elements generate every
matrix entry.
The universal fibre is generated by those
entries, so has dimension at most \(2C_0\).
In rank two only the strict linear space
occurs.
The fixed-weight topology and (LS6) give
the same bound for the completed fibre
over \(\mathfrak m_B\).

For the diagonal characters outside that
fibre, use the following finite detection
of abelian characters.
There are \(C_0+1+d\) elements \(s_i\)
in \(\ker\Psi\), and \(d\) local elements
\(s_i'\) in \(\ker\Psi\), with one fixed
exponent \(h\ne0\), such that:
\begin{enumerate}
\item a continuous character of the CM group
into a finite abelian \(\ell\)-group,
trivial on the \(s_i\), has exponent dividing \(h\);
\item if additionally
\(\chi(g^*)=\chi(g)\), its triviality on
the \(s_i'\) has the same consequence.
\end{enumerate}
To construct them, choose \(C_0+1\)
evaluations detecting
\(\operatorname{Hom}(\Gamma,k)\), by (AF4),
and append the \(d\) ordinary unit generators.
Split a character into its \(c\)-even and
\(c\)-odd parts, using that two is invertible.
The even part extends to \(\Gamma\);
the first evaluations detect it.
The odd part, after its local inertia
values vanish, factors through the finite
CM class group.
Its exponent is \(h\).
Concretely the first evaluations may be
replaced by
\((u\,cuc^{-1})^{\#k^\times}\);
they then belong to \(\ker\Psi\), without
changing their detection of an
\(\ell\)-primary character.
The \(s_i'\) are the \(d\) local unit
generators alone.
If a finite abelian \(\ell\)-group character
escaped these tests, its nonzero quotient
of order \(\ell\) would contradict the
corresponding \(k\)-linear evaluation
property.
This proves the assertion for every such
finite target, not just \(k\).

Now cut the lower-triangular locus by
the coefficient uniformizer, by
\(\chi_0(s_i)-1\), and in rank three by
\(\chi_1(s_i')-1\).
In each prime quotient containing these
equations, the diagonal characters equal
their residue characters.
Indeed their principal-unit parts have
finite exponent by the preceding tests.
In a characteristic-\(\ell\) domain there
is no nontrivial \(\ell\)-power torsion;
the remaining prime-to-\(\ell\) torsion is
uniquely fixed by its residue.
Polarization determines the final
diagonal character from the first.
The middle character in rank three
satisfies the star condition in item 2.
The ordinary product identities identify
the local characters with these same
ordered diagonal characters.
All off-diagonal monomials of weight
zero vanish on a triangular matrix.
Therefore this cut kills the entire
maximal ideal \(\mathfrak m_B\).

There are at most
\(2+C_0+(n-1)d\) cutting equations.
The height theorem and the characteristic
fibre bound prove, in the relaxed
constants of the frozen source,
\[
     \dim R_x/\mathfrak t_-
                    \le(n-1)d+3C_0+7<nd .                   \tag{SH1}
\]
On a chart with
\(\mathfrak t_+\not\subseteq x\), the
other triangular ideal is the unit ideal.

Consider next a supported local coincidence.
The actual finite-weight ideal \(I\subset B\)
has
\[
       \mathfrak m_B\subseteq\sqrt{I+\mathfrak l},
\]
where \(\mathfrak l\) is generated by the
coefficient uniformizer and ordinary
inertia augmentation on
\(\ker\varepsilon\).
Suppose \(\psi_i=\psi_j\) on
\(N_v=I_v\cap\ker\varepsilon\).
The local unit basis has \(d_v-1\)
members in \(N_v\).
In the system cutting out ordinary
augmentation omit the \(j\)-th character
on precisely those members.
Its values are supplied by the \(i\)-th
character and the coincidence equations.
Keep its value on the one remaining
generator.
Together with the coefficient uniformizer
this uses at most \(nd-d_v+2\) equations.
Their radical is maximal by the
finite-weight property.
Hence the corresponding quotient of \(B\)
has dimension at most \(nd-d_v+2\).

At an off-triangular prime the relative
transcendence degree of \(A/B\) is at most
one.
Choose nonzero entries \(u,v\) of opposite
weights.
For any other entry \(z\), a positive
power of \(z\) times a power of \(u\) or
\(v\) has weight zero.
These products lie in \(B\), and are
nonzero in the prime quotient.
Thus \(z\) is algebraic over
\(\operatorname{Frac}(B)[u]\).
Apply (LS6).
For a characteristic-\(\ell\) prime on
the coincidence locus, the relaxed
source bound is
\[
                         \dim R_x/P+d_v\le nd+5 .             \tag{SH2}
\]
Since \(d_v>5\), its dimension is
strictly less than \(nd\).
For triangular primes use (SH1).
Taking the supremum over prime quotients
proves the same strict bound for each
full coincidence ideal.

\subsection{The selected characteristic-prime curves}
Suppose \(P\subset R_x\) contains the
supported ideal and has
\(\dim R_x/P\ge nd+1\).
Apply the one-dimensional prime-selection
lemma already proved, with threshold \(nd\),
the coefficient uniformizer, and the
finitely many small ideals in (SH1)--(SH2).
It gives an actual prime \(p\supseteq P\)
satisfying
\[
 \ell\in p,\quad \dim R_x/p=1,\quad
 \mathfrak t_\pm\not\subseteq p,\quad
       \psi_i/\psi_j\ne1\text{ on every }N_v.                 \tag{SH3}
\]
These are every field of the source
\srcid{Nice} predicate.

The normalization of the complete
one-dimensional local domain \(R_x/p\)
is finite.
It is a complete local DVR with finite
residue field \(K\).
Since it has equal characteristic \(\ell\)
and \(K\) is perfect, it contains its unique
coefficient field lifting \(K\): its elements
are the roots of \(X^{|K|}-X\).
Successive expansion in a uniformizer identifies
this DVR with \(K[[T]]\).
Thus it gives the actual \srcid{W.Point}
coefficient ring, with kernel \(p\).
The local character product identities
give its ordinary flags by the separated
eigenvalue argument of (LS1).
The characteristic-\(\ell\) regular-coset
proof from these separated points supplies
the previously established \(N1,N3\);
the distinguished scalar Frobenius
supplies \(N4\).
Thus the selected curves give actual
nice points in the same determinant
space, rather than hypothetical points
of an unrelated ring.

\subsection{Connectedness after a determinantal section}
We use and prove the needed complete-local
connectedness statement.
If \(C\) is an excellent complete local
domain of dimension \(q\), and \(h\)
elements of its maximal ideal generate
\(J_0\), every minimal prime of \(J_0\)
has quotient dimension at least \(q-h\),
and the graph of these minimal primes
with edges
\[
                        \dim C/(P+Q)\ge q-h-1                \tag{SH4}
\]
is connected.
The component bound is the height theorem
and catenarity.

For the graph, normalize \(C\), obtaining
a finite complete normal local domain.
A disconnected graph would partition its
components after the same \(h\) equations
into two nonempty unions with intersection
dimension less than \(q-h-1\).
This reduction is legitimate under the
finite normalization map: dimensions are
preserved, each minimal component upstairs
has dimension at least \(q-h\), and hence
none maps entirely into that smaller
intersection.
Regroup the minimal components upstairs
according to the two unions downstairs.

Localize at a prime minimal over the
intersection of the two unions.
Its height is at least \(h+2\).
The localization is excellent and normal,
so its completion is a domain of that
dimension \(q'\ge h+2\).
The two unions give a disconnection of
the punctured spectrum after the same
\(h\) equations.
Let their ideals be \(J_1,J_2\).
Neither is maximal-primary and
\(\sqrt{J_1+J_2}=\mathfrak m\),
\(\sqrt{J_1\cap J_2}=\sqrt{J_0}\).
The Cech complex on \(h\) generators
gives \(H^{q'-1}_{J_0}=0\).
The
\href{https://stacks.math.columbia.edu/tag/0EB0}
{Hartshorne--Lichtenbaum theorem}
gives \(H^{q'}_{J_1}=H^{q'}_{J_2}=0\).
Mayer--Vietoris would give
\(H^{q'}_{\mathfrak m}=0\), contrary
to top local-cohomology nonvanishing.
This proves (SH4), including the case
of a nonequidimensional quotient.

Put \(D_0=3d+1\), and consider the even/odd
presentation (LS11), with \(G=My\).
Introduce independent variables
\(Z_{ij},T_j\) over its regular ambient
power-series ring.
By the explicit incidence parametrization
already proved, the radical of
\[
                               (ZT,\det Z)
\]
is prime of height at most \(g\).
Its quotient domain has dimension at least
\(D_0+1+b+g^2+g\).
Cut by \(Z_{ij}-M_{ij}\), by \(T_j-y_j\),
and by the \(b\) even relations.
All \(g^2+g+b\) equations belong to the
maximal ideal.
Statement (SH4) shows that
\[
             I^\sharp=(\det M)\subset R_x
            \text{ has its minimal-prime graph at threshold }D_0:
                                                               \tag{SH5}
\]
every minimal prime has quotient dimension
at least \(D_0+1\), and adjacent ones may
be joined with
\(\dim R_x/(P+Q)\ge D_0\).
Radicals and the surjection from the
ambient ring preserve this graph and
its dimensions.

In the actual exceptional chart \(g>0\).
Otherwise (LS11) would give
\(\mathfrak eR_x=0\).
Faithful flatness after localizing
\(A\) at \(x\) then gives a single
\(s\notin x\) annihilating the finitely
generated ideal \(\mathfrak e\subset A\).
The actual polynomial families after
(LS4) prove that the nonexceptional
locus is dense in the characteristic
core.
It therefore meets its basic
neighborhood \(D(s)\).
But \(s\mathfrak e=0\) makes every
prime in \(D(s)\) contain \(\mathfrak e\),
a contradiction.

The adjugate identity gives
\((\det M)y_j=0\) in \(R_x\).
Every prime avoiding
\(\mathfrak e=(y_j)\) consequently
contains \(I^\sharp\).
Every minimal component in (SH5)
is off both triangular loci, by
its dimension lower bound and (SH1).
If it avoids \(\mathfrak e\), it is
Type I.

\subsection{The actual thick representation on an exceptional component}
Let \(P\) be a minimal prime in (SH5)
containing \(\mathfrak e\).
Over \(R_x/P\), the singular matrix \(M\)
has a nonzero kernel vector \(\nu\).
Choose it over the fraction field and
clear its finitely many denominators.
The operator
\[
          D(f)=\sum_j\nu_j\frac{\partial f}{\partial y_j}
                                               \pmod P        \tag{SH6}
\]
kills the even relations: their odd
derivatives belong to \((y_j)\).
It kills \(My\) because \(M\nu=0\).
Thus it is a derivation \(R_x\to R_x/P\),
zero on the even subring and on \(B\),
and nonzero on at least one \(y_j\).

Define actual matrices by
\[
 s_D(g)=
 \begin{pmatrix}
 \rho_{00}(g)&D\rho_{01}(g)&\rho_{02}(g)\\
 0&\rho_{11}(g)&0\\
 \rho_{20}(g)&D\rho_{21}(g)&\rho_{22}(g)
 \end{pmatrix}
                    \quad\text{over }R_x/P .                 \tag{SH7}
\]
Let \(\pi:R_x\to R_x/P\) be the quotient map.
For a matrix \(a\), require that its four odd
entries vanish under \(\pi\) and that the
derivatives of its five even entries vanish.
The matrices with these properties are closed
under multiplication.
For two such matrices, the Leibniz rule gives
\[
 \operatorname{th}(ab)
       =\operatorname{th}(a)\operatorname{th}(b),
 \qquad
 \operatorname{th}(a-cI)
       =\operatorname{th}(a)-\pi(c)I
                          \quad\text{if }D(c)=0,
\]
where \(\operatorname{th}\) retains the even
entries, differentiates entries \((0,1),(2,1)\),
and sets entries \((1,0),(1,2)\) to zero.
In particular \(\operatorname{th}(I)=I\), and
\(s_D=\operatorname{th}\circ\rho\) is a group
representation.

The original polarization and its ordered
ordinary identity imply
\[
 \prod_{i=0}^{2}
   \bigl(\rho((\operatorname{loc}_v g_i)^*)
                    -\psi_{v,2-i}(g_i)I\bigr)=0.             \tag{SH8}
\]
Indeed apply \(a\mapsto J\,{}^taJ\) to the original
identity with the arguments in reversed order;
transpose reverses the product.
Every factor in the ordinary identity and
(SH8) has the two properties just specified,
since \(D\psi_{vi}=0\).
Applying the multiplicative map
\(\operatorname{th}\) therefore proves exactly
\[
 \begin{split}
 \prod_{i=0}^{2}
   \bigl(s_D(\operatorname{loc}_v g_i)
                    -\pi(\psi_{vi}(g_i))I\bigr)&=0,\\
 \prod_{i=0}^{2}
   \bigl(s_D((\operatorname{loc}_v g_i)^*)
                    -\pi(\psi_{v,2-i}(g_i))I\bigr)&=0.
 \end{split}                                                \tag{SH9}
\]
These are the two ordered local identities in
\srcid{IsThickMod}.
Polarization is used here for the original
\(\rho\); \srcid{IsThickMod} requires the
identities (SH9) for \(s_D\).
Continuity is retained because a
derivation sends
\(\mathfrak m_{R_x}^{N+1}\) into
\(\mathfrak m_{R_x/P}^N\).

This middle-character extension is
nonsplit after every injective embedding
of \(R_x/P\) into a field.
A splitting vector may be normalized
by \(v_1=1\).
Subtract its eigenvector equations
at \(g_0\) and \(g_0^*\).
The marked normalization makes the
off-diagonal entries equal.
Their outer diagonal difference is
a unit and their middle difference
is zero.
Therefore \(v_0=v_2=0\).
The remaining eigenvector equations
give
\(D\rho_{01}=D\rho_{21}=0\).
Polarization gives the other odd
derivatives zero, and the even
derivatives already vanish.
The ideal of odd entries is
\((y_j)\); the derivation rule would
make \(D(y_j)=0\) for every \(j\),
contrary to (SH6).
This proves all parts of the actual
\srcid{IsThickMod} predicate:
the same even part, the specified
zero middle row, both ordered local
identities, continuity, and
nonsplitting after scalar extension.

Off the exceptional chart use
\(I^\sharp=0\).
The ordinary presentation (LS10)
is \(b\) equations in \(b+D_0\)
variables, so (SH4) gives its graph
at the same threshold.
All minimal components are off
both triangular loci by (SH1),
and none contains \(\mathfrak e\),
which is the unit ideal on this
chart.
They are all Type I.

Thus the specialized sharp ideal
has actually been constructed.
It is zero off the exceptional
chart, lies in every Type-I prime,
has the required connected
minimal-prime graph, and every
minimal component is Type I or
has the explicit thick model
(SH7).

\paragraph{Exact source correspondence.}
The small-fibre generators are
\srcid{FinrankStructMapYRange},
\srcid{FinrankStructMapSRange},
\srcid{FinrankFrameSpaceEq},
and \srcid{FibDim}.
Abelian detection and the two small
loci are
\srcid{AbEval.ab\_gen},
\srcid{CharpolyCoeffsInSubalgebra.l51},
\srcid{TameInertiaCondition.l52},
and
\srcid{KrullDimLtOfNotLeTriangular}.
The curve producer is
\srcid{CurveDVRSpecialization.%
ExistsNicePrimeAboveOfLargeDim}.
The determinant section and thick model
expand \srcid{Detsec.detsec},
\srcid{Oddm.oddm},
\srcid{AnnOfFlat.gpos},
\srcid{ThickOfDer.thick}, and
\srcid{ChainedOfComap.sharp3}.
Except for the indicated curve module,
these paths are under
\srcid{FrobeniusSmoothVectors}, within
\srcid{Development.AutomorphicArithmetic.AutomorphicVarieties}
at \srcid{alpoge-source-7cbf7a0c}.

These producers supply the actual
flat-chart, sharp-ideal, dichotomy
and curve-selection hypotheses in
the already proved local/global
propagation argument.
The remaining insertion into its
supported locus uses the common
Hecke module and its completed
support inclusions; it is not a
consequence of connectedness alone.
